\documentclass[11pt]{article}
\usepackage[T1]{fontenc}
\usepackage{amsmath,amssymb,amsthm}
\usepackage[margin=0.88in]{geometry}
\usepackage{booktabs}
\usepackage{array}
\usepackage{hyperref}
\usepackage{microtype}

\hypersetup{
  pdftitle={Exact Diameters and Sorting Times for Full Cycles of Orders 6, 7, and 8},
  pdfauthor={Xiaojian Zeng},
  pdfsubject={Exact exhaustive computation in finite Schreier graphs}
}

\newtheorem{theorem}{Theorem}
\newtheorem{lemma}[theorem]{Lemma}
\newcommand{\Sort}{\operatorname{Sort}}
\newcommand{\dist}{\operatorname{dist}}
\newcommand{\cyc}[1]{\langle #1\rangle}

\title{Exact Diameters and Sorting Times for Full Cycles of Orders 6, 7, and 8}
\author{Xiaojian Zeng\\
PBC School of Finance, Tsinghua University, 100190\\
\texttt{xiaojianzeng@foxmail.com}}
\date{}

\begin{document}
\maketitle

\begin{abstract}
Let $\Gamma_n$ have the full $n$-cycles as vertices, with an edge from $c$ to
$s_i c s_i$ for each adjacent transposition $s_i=(i,i+1)$.  We determine by
complete exact enumeration
\[
 (D_6,\Sort_6)=(7,6),\qquad (D_7,\Sort_7)=(10,10),\qquad
 (D_8,\Sort_8)=(14,13).
\]
Two separately written programs use canonical cycle words and one-line
permutation maps, respectively.  Each constructs the graph independently and
runs breadth-first search from every vertex; the resulting distance
histograms and deterministic hashes agree.  Explicit shortest paths are
included.  The sorting values reproduce computations reported by
Adin--Bagno--Roichman; no priority claim is made for the diameter values, and
the calculation has no implication for their asymptotic or unimodality
conjectures.
\end{abstract}

\section{Statement and context}

Let $C_n$ be the conjugacy class of full cycles in $S_n$, let
$s_i=(i,i+1)$ for $1\leq i<n$, and let $\Gamma_n$ be the simple undirected
graph with edges $\{c,s_i c s_i\}$.  Put $c_n=(1,2,\ldots,n)$ and
\[
 D_n=\max_{u,v\in C_n}\dist_{\Gamma_n}(u,v),\qquad
 \Sort_n=\max_{v\in C_n}\dist_{\Gamma_n}(c_n,v).
\]
These are the diameter and Coxeter cyclic sorting time introduced by
Adin, Bagno, and Roichman~\cite{AdinBagnoRoichman2026}.  They identify
$\Gamma_n$ with a Schreier graph and give general bounds.  They also report
$D_5=5>4=\Sort_5$ and verify their proposed sorting formula through $n=12$;
in particular, their computation already gives the three sorting values below.

\begin{theorem}\label{thm:main}
For the graph just defined,
\[
\begin{array}{c|ccc}
n&6&7&8\\ \hline
\Sort_n&6&10&13\\
D_n&7&10&14
\end{array}
\]
The maxima range over all $(n-1)!$ vertices, and the diameter maxima range over
all ordered pairs of vertices.
\end{theorem}

The contribution is the exact finite determination requested for orders
$6,7,8$, including reproducible diameter computations, complete distance
distributions, and shortest witnesses.  The certificate architecture is close
in spirit to other finite graph certificates that reconstruct the instance and
check the decisive predicate in separately written exact implementations; see,
for example, Zeng et al.~\cite[Sections 6--7]{ZengEtAl2026FiniteCertificates}.
That comparison is methodological: the mathematical problems are unrelated.

\section{Exhaustive certificate}

A cycle is represented by its unique rotation
\[
  (1,a_2,\ldots,a_n),
\]
where the displayed permutation maps each entry to its successor and the last
entry to $1$.  Reversal is not identified.  We abbreviate this word as
$\cyc{a_2\cdots a_n}$ below.  Thus every permutation of $2,\ldots,n$ occurs
once, giving exactly $(n-1)!$ vertices.

\begin{lemma}\label{lem:edge}
In the canonical cycle-word representation, $s_i c s_i$ is obtained by
interchanging the labels $i$ and $i+1$ everywhere in the word and then rotating
the result to begin with $1$.
\end{lemma}

\begin{proof}
Conjugation relabels every entry $x$ of a cycle as $s_i(x)$.  A cyclic rotation
does not change the resulting cycle, so rotating to put $1$ first gives its
unique canonical word.
\end{proof}

Verifier A enumerates the canonical words lexicographically and applies
Lemma~\ref{lem:edge}.  Verifier B instead enumerates all one-line maps in
$S_n$, retaining a map $p$ precisely when the orbit of $1$ first returns after
$n$ steps.  It constructs a neighbor by literal map composition
\[
 q(x)=s_i\bigl(p(s_i(x))\bigr),
\]
without using A's word-swap routine.  Both programs check connectivity,
involutivity of each labeled edge, absence of loops and duplicate generator
neighbors, and the following counts.

\begin{center}
\begin{tabular}{crrr}
\toprule
$n$ & vertices & degree & undirected edges\\
\midrule
6&120&5&300\\
7&720&6&2,160\\
8&5,040&7&17,640\\
\bottomrule
\end{tabular}
\end{center}

Each verifier performs ordinary unweighted breadth-first search from every
vertex.  Hence it computes every ordered-pair distance: the maximum in the row
rooted at $c_n$ is $\Sort_n$, while the maximum over all rows is $D_n$.  The
complete rooted histograms $H_n^{\rm root}(d)$ are

\begin{center}\small
\begin{tabular}{c|rrrrrrrrrrrrrrr}
\toprule
$n\backslash d$&0&1&2&3&4&5&6&7&8&9&10&11&12&13\\
\midrule
6&1&5&14&27&36&28&9&&&&&&&\\
7&1&6&20&48&91&140&169&146&78&20&1&&&\\
8&1&7&27&76&172&329&544&784&970&978&729&344&77&2\\
\bottomrule
\end{tabular}
\end{center}
and the all-ordered-pairs histograms $H_n^{\rm all}(d)$ are

\begin{center}\scriptsize
\begin{tabular}{c|rrrrrrrr}
\toprule
$n\backslash d$&0&1&2&3&4&5&6&7\\
\midrule
6&120&600&1,620&3,134&4,268&3,424&1,162&72\\
7&720&4,320&14,400&34,962&67,284&103,662&123,828&103,614\\
8&5,040&35,280&135,936&381,656&866,404&1,676,632&2,826,516&4,131,396\\
\bottomrule
\end{tabular}

\medskip
\begin{tabular}{c|rrrrrrr}
\toprule
$n\backslash d$&8&9&10&11&12&13&14\\
\midrule
6&&&&&&&\\
7&51,852&12,642&1,116&&&&\\
8&5,081,316&4,973,336&3,461,236&1,486,704&319,284&20,784&80\\
\bottomrule
\end{tabular}
\end{center}
The rows sum to $(n-1)!$ and $((n-1)!)^2$, respectively.  Their last nonzero
levels prove the upper bounds in Theorem~\ref{thm:main}; the paths in
Appendix~\ref{app:paths} attain them and prove the matching lower bounds.

For deterministic comparison, vertices are serialized in lexicographic order,
undirected edges by sorted endpoint pairs, and distance rows in vertex order.
The two implementations agree on all three SHA-256 streams:

\begin{center}\scriptsize\ttfamily
\begin{tabular}{ccl}
\toprule
$n$&stream&SHA-256\\
\midrule
6&vertices&80515474c51c1389bbf0de1fb7efe3eaefa5db03b42b3d33b8073f91249478c4\\
7&vertices&2c2f4bf673ff0ad5f4cc54c54337813877c4f84307ae76bcceac38aa62e24147\\
8&vertices&81c19e18a046ee68d22e66a3ad83cd7996604685829d84bf947f5320d33a8b11\\
6&edges&e8000a7ce9605187753a9a742d888d71e7182f77a9c74fd6615cc8acdf7cd66c\\
7&edges&a3c33fff0b1e78d8287ceff39f1464571d9d541786ddb86ee29d1f87b9ddb0a2\\
8&edges&33bc18caa3367881413fa8f20e21dfe8936342291caf41187ce37f5c2405c83f\\
6&distances&526c7ef0cd53c0d83e42e0f63249f7496b9c030d1109dfa683ac21303245b1f7\\
7&distances&c62e059528260bca71a36fa5a0ba4e118850e03378cc5a9dd8b6b6a6b669a081\\
8&distances&4c1bc1c3e6aeeff3e22198a39dd09fe86841484c08ad5ef8183410ac26073fb6\\
\bottomrule
\end{tabular}
\end{center}

Both submitted runs used CPython 3.10.12 on Linux and only the Python standard
library.  From the project root they are replayed by
\begin{verbatim}
for n in 6 7 8; do
  python3 evidence/verifier_a.py --n "$n" \
    --output "evidence/verifier_a_n${n}.json"
  python3 evidence/verifier_b.py --n "$n" \
    --compare-a "evidence/verifier_a_n${n}.json" \
    --output "evidence/verifier_b_n${n}.json"
done
\end{verbatim}
Verifier B also replays every path produced by A using literal map composition
and confirms that its length is the independently computed endpoint distance.

\section{Scope and disclosure}

This is a finite exhaustive theorem for $n\in\{6,7,8\}$.  It neither proves
nor disproves the asymptotic Conjecture 7.1 or the unimodality Conjecture 7.2 of
Adin--Bagno--Roichman~\cite{AdinBagnoRoichman2026}.  The three sorting values
reproduce their finite experiments.  No assertion of priority or journal
acceptance is made for the three diameter values.

OpenAI Codex assisted with literature searching, implementation and audit
organization, and initial manuscript preparation.  Language-model output is
not mathematical evidence.  The result rests on the finite argument above,
the released exact programs and records, and independently replayable paths.
No proof assistant was used.

\appendix
\section{Explicit shortest paths}\label{app:paths}

An arrow labeled $i$ means conjugation by $s_i=(i,i+1)$.  The abbreviation
$\cyc{a_2\cdots a_n}$ denotes $(1,a_2,\ldots,a_n)$.  The displayed generator
lists are recorded as an additional compact replay check.

\paragraph{$n=6$, sorting witness.}
The generator list is $(4,3,4,2,3,4)$, and
\[
\begin{aligned}
\cyc{23456}&\xrightarrow{4}\cyc{23546}\xrightarrow{3}\cyc{24536}
\xrightarrow{4}\cyc{25436}\xrightarrow{2}\cyc{35426}\\
&\xrightarrow{3}\cyc{45326}\xrightarrow{4}\cyc{54326}.
\end{aligned}
\]

\paragraph{$n=6$, diameter witness.}
The generator list is $(4,5,4,3,2,3,4)$, and
\[
\begin{aligned}
\cyc{23564}&\xrightarrow{4}\cyc{23465}\xrightarrow{5}\cyc{23456}
\xrightarrow{4}\cyc{23546}\xrightarrow{3}\cyc{24536}\\
&\xrightarrow{2}\cyc{34526}\xrightarrow{3}\cyc{43526}
\xrightarrow{4}\cyc{53426}.
\end{aligned}
\]

\paragraph{$n=7$, sorting and diameter witness.}
Here the sorting path is also a diameter path.  The generator list is
$(6,5,6,3,4,5,6,2,1,2)$, and
\[
\begin{aligned}
\cyc{234567}&\xrightarrow{6}\cyc{234576}\xrightarrow{5}\cyc{234675}
\xrightarrow{6}\cyc{234765}\xrightarrow{3}\cyc{243765}\\
&\xrightarrow{4}\cyc{253764}\xrightarrow{5}\cyc{263754}
\xrightarrow{6}\cyc{273654}\xrightarrow{2}\cyc{372654}\\
&\xrightarrow{1}\cyc{654237}\xrightarrow{2}\cyc{654327}.
\end{aligned}
\]

\paragraph{$n=8$, sorting witness.}
The generator list is $(7,5,6,7,4,3,4,2,3,4,1,2,3)$, and
\[
\begin{aligned}
\cyc{2345678}&\xrightarrow{7}\cyc{2345687}\xrightarrow{5}\cyc{2346587}
\xrightarrow{6}\cyc{2347586}\xrightarrow{7}\cyc{2348576}\\
&\xrightarrow{4}\cyc{2358476}\xrightarrow{3}\cyc{2458376}
\xrightarrow{4}\cyc{2548376}\xrightarrow{2}\cyc{3548276}\\
&\xrightarrow{3}\cyc{4538276}\xrightarrow{4}\cyc{5438276}
\xrightarrow{1}\cyc{7625438}\xrightarrow{2}\cyc{7635428}\\
&\xrightarrow{3}\cyc{7645328}.
\end{aligned}
\]

\paragraph{$n=8$, diameter witness.}
The generator list is $(6,5,6,4,5,3,4,5,2,3,4,5,6,7)$, and
\[
\begin{aligned}
\cyc{2347856}&\xrightarrow{6}\cyc{2346857}\xrightarrow{5}\cyc{2345867}
\xrightarrow{6}\cyc{2345876}\xrightarrow{4}\cyc{2354876}\\
&\xrightarrow{5}\cyc{2364875}\xrightarrow{3}\cyc{2463875}
\xrightarrow{4}\cyc{2563874}\xrightarrow{5}\cyc{2653874}\\
&\xrightarrow{2}\cyc{3652874}\xrightarrow{3}\cyc{4652873}
\xrightarrow{4}\cyc{5642873}\xrightarrow{5}\cyc{6542873}\\
&\xrightarrow{6}\cyc{7542863}\xrightarrow{7}\cyc{8542763}.
\end{aligned}
\]

\bibliographystyle{plain}
\bibliography{references}

\end{document}
